% Post-CP18 source. Four complete actual-file and primary-source audits passed; fixed line bundles only.
\subsection{Finite polynomial Hom complexes for fixed determinant presentations}

\begin{proposition}[Proved: all-coefficient finite Hom and diagram models]
\label{D:fixed-determinant-finite-Hom}
Let $D$ be a projective Cohen--Macaulay scheme of pure dimension
one over an algebraically closed field $k$. Put
$b_D=h^0(D,\mathcal O_D)+h^1(D,\mathcal O_D)$.
Fix globally generated line bundles $L_i$ and positive integers
$r_i$, and put $d_i=\deg_D L_i$. On the open set of generating
tuples in $H^0(D,L_i)^{r_i+1}$ consider the universal sequence
\[
 0\longrightarrow L_i^{-1}
   \xrightarrow{s_i}\mathcal O_D^{r_i+1}
   \longrightarrow V_i\longrightarrow0.
                                                        \tag{D.FH.1}
\]
All constant sheaves in this formula are pulled back to the
coefficient scheme. There are finite graded $k$-vector spaces
$H_{ij}$ and differentials $\Delta_{ij}(s)$ on
$H_{ij}\otimes_k T$, for every coefficient $k$-algebra $T$,
with the following properties.
\begin{enumerate}
\item The construction is given by fixed polynomial formulas
in the coordinates of the tuples $s_i$, commutes with every
homomorphism of coefficient algebras, and computes
\[
 (H_{ij}\otimes_k T,\Delta_{ij})
       \simeq R\operatorname{Hom}_{D_T}(V_i,V_j).
                                                        \tag{D.FH.2}
\]
The degrees of the entries of $\Delta_{ij}$ are at most two.
No flatness or base-change assertion for the cohomology of
this complex is assumed.
\item Its total dimension is at most
\[
 \begin{split}
 B_{ij}={}&b_D(r_i+1)(r_j+1)
 +(b_D+d_j)(r_i+1)\\
 &+(b_D+d_i)(r_j+1)+b_D+d_i+d_j.
 \end{split}                                           \tag{D.FH.3}
\]
It is concentrated in total degrees $-1,0,1,2$.
\item Actual morphisms $V_i\to V_j$ over $T$ are represented
by degree-zero cycles in this finite complex. Equality of
two representatives is expressed by one degree-minus-one
boundary witness. Composition is represented by a polynomial
map of total degree at most five in the object and arrow
coordinates. Inverse witnesses therefore give an exact
polynomial presentation of actual bundle isomorphisms over
every $T$, including nonreduced algebras.
\end{enumerate}

In particular, suppose $b_D\le cn$, $d_i\le cnr_i$, and
$1\le r_i\le r$. For a diagram with $O(f)$ vertices, arrows,
and relations, whose relation lengths are bounded by a fixed
integer, its actual bundles, arrows, and relations admit
these presentations using $O(fnr^2)$ coordinates and scalar
equations of bounded degree. The constants depend on $c$
and the fixed diagram pattern, not on $n,r,f$.
The generating-tuple conditions are open conditions.
The line bundles $L_i$ and the curve $D$ are fixed background
data; this proposition does not parameterize their variation.
\end{proposition}

\begin{proof}
\emph{The line-bundle cohomology budget.}
A general section of a globally generated line bundle on
$D$ avoids its finitely many associated points. It is a
nonzerodivisor, with a zero scheme of length equal to the
degree of the line bundle. The sequences defined by such
sections show that
\[
 \begin{split}
 h^0(L_i)+h^1(L_i)&\le b_D+d_i,\\
 h^0(L_i^{-1})+h^1(L_i^{-1})&\le b_D+d_i,\\
 h^0(L_iL_j^{-1})+h^1(L_iL_j^{-1})
                                      &\le b_D+d_i+d_j.
 \end{split}                                           \tag{D.FH.4}
\]
For the last inequality, tensor the section sequence for
$L_i$ by $L_j^{-1}$, and use the preceding bound for
$L_j^{-1}$. These arguments use the full nonreduced curve;
the Cohen--Macaulay hypothesis rules out embedded associated
points.

\emph{A fixed two-column Cech model.}
Embed $D$ into projective space. Choose two hyperplanes whose
intersection misses $D$ and neither of which contains an
irreducible component. Their complements give an affine
cover $D=U_0\cup U_1$ with affine intersection. Thus the
ordered Cech complex of any coherent sheaf has just the
two terms in Cech degrees zero and one, and computes its
cohomology. The exact primary input is
\citep[Lemmas~30.2.5--30.2.6 and the complete proofs;
Section~30.5, Lemmas~30.5.1--30.5.2]{C_StacksFiniteCechHom}.
The complexes on $D_T$ are obtained by tensoring the fixed
component Cech complexes over $k$ with $T$.

For each of the finitely many line bundles occurring below,
choose, once over $k$, a splitting of its two-term Cech
complex onto its finite-dimensional cohomology. Write its
maps as $\iota,\rho,h$, with
\[
 \rho\iota=1,\qquad
 \iota\rho=1-dh-hd,\qquad
 \rho h=h\iota=h^2=0.                                 \tag{D.FH.5}
\]
Here $d$ is the Cech differential, $h$ lowers Cech degree,
and the target of the splitting has zero differential.
Explicitly, split the kernel in degree zero and split the
image in degree one; invert $d$ between the two complementary
subspaces to define $h$. No finite-dimensionality of the
original spaces of affine sections is required. The maps
are fixed $k$-linear maps, and (D.FH.5) remains valid after
tensoring with every $T$.

Put $P_i=[L_i^{-1}\xrightarrow{s_i}\mathcal O_D^{r_i+1}]$
in internal degrees $-1,0$. Form the total Cech Hom complex
$C_{ij}$ for $\mathcal Hom^\bullet(P_i,P_j)$. Its differential
has the form
\[
                         d+\delta,
\]
with the usual total-complex signs. The component $\delta$
is the internal Hom differential. It raises internal degree,
preserves Cech degree, and is linear in the tuple coordinates.
In particular
\[
 d\delta+\delta d=0,\qquad \delta^2=0,\qquad
                         h\delta h=0.                  \tag{D.FH.6}
\]
The last equality follows because $h$ sends Cech degree one
to zero and is zero on degree zero. Use the direct sums of
the fixed splittings (D.FH.5) in all internal degrees.

\emph{An explicit finite transfer, without variable bases.}
On the finite graded cohomology space define
\[
 \begin{split}
 I&=\iota-h\delta\iota,\qquad
 P=\rho-\rho\delta h,\\
 \Delta&=\rho\delta\iota-\rho\delta h\delta\iota.
 \end{split}                                           \tag{D.FH.7}
\]
The signs correspond to the convention in (D.FH.5).
Substitution using (D.FH.5)--(D.FH.6) gives
\[
 (d+\delta)I=I\Delta,\quad
 P(d+\delta)=\Delta P,\quad PI=1,\quad
 IP=1-(d+\delta)h-h(d+\delta).                          \tag{D.FH.8}
\]
For example the terms beyond (D.FH.7) vanish because
$h\delta h=0$; the cross term
$h\delta\iota\rho\delta h$ also vanishes after replacing
$\iota\rho$ by $1-dh-hd$ and using (D.FH.6).
These identities prove $\Delta^2=0$ and a chain homotopy
equivalence. They hold polynomially over every coefficient
ring, not only at field-valued points. The differential
has degree at most two, while $I$ and $P$ have degree at most
one in the object coordinates.

The finite space here is the direct sum of the cohomology
of the internal Hom terms
\[
 L_j^{-1\,\oplus(r_i+1)},\quad
 \mathcal O_D^{\oplus(r_i+1)(r_j+1)}\oplus L_iL_j^{-1},
 \quad L_i^{\oplus(r_j+1)},                             \tag{D.FH.9}
\]
in internal degrees $-1,0,1$, respectively. In the first
term, the notation means $(L_j^{-1})^{\oplus(r_i+1)}$.
Equations (D.FH.4) give precisely (D.FH.3).
On the open generating locus, $P_i$ is a locally free
resolution of the vector bundle $V_i$. The Cech Hom complex
therefore computes the derived Hom in (D.FH.2), over every
$T$. Exactness of the presentation is preserved by every
coefficient change because its quotient is locally free.

\emph{Composition and equality witnesses.}
The two-column ordered Cech model has its usual associative
composition product, with the total-complex signs. For
cochains in Cech degrees zero and one, the products use
the first or last restriction respectively; a product of
two degree-one Cech cochains is zero. For degree-zero
cycles $a\in H_{ij}\otimes T$ and
$b\in H_{jk}\otimes T$, define
\[
             m_2(b,a)=P_{ik}\bigl(I_{jk}(b)\,I_{ij}(a)\bigr).
                                                        \tag{D.FH.10}
\]
It is a cycle and represents actual composition. Its total
degree is at most five: two arrow coordinates, at most one
object factor from each inclusion, and at most one from the
projection. This product need not be strictly associative
on all finite cochains, and no such assertion is needed;
on degree-zero cohomology it is the associative bundle
composition. A path of any fixed length can instead be
represented by multiplying all its lifted Cech cycles first
and then applying the final projection. Its polynomial
degree is bounded by twice the path length plus one.

Every actual arrow is represented by a cycle, since
$H^0$ of (D.FH.2) is its ordinary Hom module. A proposed
equality between paths is exact precisely when the
difference of their representatives is $\Delta w$ for
some $w$ in degree $-1$. This is a finite list of polynomial
equations with a boundary witness. In particular two arrows
are mutual inverses exactly when their two compositions
differ from the identity cycles by such boundaries.
The identity cycles may be obtained by applying $P$ to the
fixed identity cochains. Thus every assertion about arrows
has its stated all-coefficient meaning; no passage from
fibrewise invertibility to an unproved lifting statement
has been made.

Finally, the object tuple uses
$(r_i+1)h^0(L_i)=O(nr_i^2)$ coordinates. Formula (D.FH.3)
bounds the cycles, boundary witnesses, and their scalar
equations by $O(nr^2)$ for each of the fixed local diagram
patterns. A bounded-length relation needs only its source,
target, and constituent Hom complexes, hence a uniformly
bounded number of them. Summing over $O(f)$ patterns gives
the claimed budget. All splittings, bases, and the finitely
many resulting multiplication coefficients are chosen once
for the fixed line bundles. They are background constants,
not additional matrix-valued moduli coordinates.
\end{proof}

\begin{remark}[What this finite model does and does not provide]
The same construction applied to $P_i$ itself gives a finite
complex of total dimension $O(nr_i)$ computing the cohomology
of $V_i$ over every coefficient algebra; its differential
is linear, since there is only one internal step. Fixed
restriction and tensor diagrams can be treated by the same
finite-width construction. This represents the derived
complexes, but does not assert cohomological flatness or
bound the equations for a chosen cohomology-rank stratum.
Likewise the proposition does not choose moving determinant
lines, realize arbitrary finite-cycle data as full modules,
or lift an arbitrarily prescribed cycle arrow to the
resolution. Those are distinct conditions in an eventual
hereditary unit-object atlas. No representation count or
finite-generation assertion follows from the proposition
alone.
\end{remark}
